% Generated draft: requires mathematical and source review.
\documentclass[11pt]{article}
\usepackage[margin=25mm]{geometry}
\usepackage{fontspec}
\setmainfont{Latin Modern Roman}
\usepackage{amsmath,amssymb,mathtools,unicode-math}
\setmathfont{Latin Modern Math}
\usepackage{enumitem,needspace,xurl,hyperref}
\setlength{\parindent}{0pt}
\setlength{\parskip}{5pt}
\setlength{\emergencystretch}{3em}
\newcommand{\N}{\mathbb{N}}
\newcommand{\Z}{\mathbb{Z}}
\newcommand{\Q}{\mathbb{Q}}
\newcommand{\R}{\mathbb{R}}
\newcommand{\C}{\mathbb{C}}
\newcommand{\F}{\mathbb{F}}
\newcommand{\A}{\mathbb{A}}
\newcommand{\PP}{\mathbb{P}}
\newcommand{\E}{\mathbb{E}}
\newcommand{\eps}{\varepsilon}
\newcommand{\dd}{\,\mathrm{d}}
\begin{document}
\section*{UnsolvedMath ID: 11000015}
\textbf{Source problem code:} catalogue entry
\subsection{Definitions and mathematical statement}
For each \(g\geq2\), let \(\mathcal I_g\) be the Torelli group in the source's surface convention. In the usual closed oriented case,
\[
 \mathcal I_g=\ker\bigl(\operatorname{Mod}_g\to
 \operatorname{Sp}(H_1(S_g;\Z))\bigr).
\]
Let \(\mathfrak t_g\) denote its Malcev Lie algebra, conventionally over \(\Q\). This is a pronilpotent Lie algebra associated to the group's rational unipotent completion. Retain the source's grading and \(\operatorname{Sp}\)-action; the excerpt does not specify the coefficient convention, filtration, or category of representations.

\emph{Problem 2.12.} Determine, for every \(g\geq2\), the irreducible factors of every graded piece of \(\mathfrak t_g\) as \(\operatorname{Sp}\)-modules. If the intended grading is the lower-central grading, with \(L^1\mathfrak t_g=\mathfrak t_g\) and \(L^{d+1}\mathfrak t_g=[\mathfrak t_g,L^d\mathfrak t_g]\) in the appropriate completed sense, these pieces are
\[
 \operatorname{gr}^L_d\mathfrak t_g
 =L^d\mathfrak t_g/L^{d+1}\mathfrak t_g,\qquad d\geq1.
\]
In a finite-dimensional semisimple rational-representation setting, the requested data can be written as the irreducible types and multiplicities in
\[
 \operatorname{gr}_d\mathfrak t_g\cong\bigoplus_\lambda
 V_\lambda^{\oplus m_\lambda(g,d)},
 \qquad\operatorname{Sp}=\operatorname{Sp}(H_1(S_g;\Q)).
\]
This notation is conditional: the source does not establish finite-dimensionality or complete reducibility in every genus, especially the included genus \(2\), or identify lower-central and any weight grading. Those conditions must not be assumed merely to force a finite decomposition.

The source says Hain [Ha3] gives an explicit, reasonably simple presentation of \(\mathfrak t_g\) when \(g>6\), while the representation problem still appears involved. The strict range \(g>6\) is retained. Its trailing discussion credits Chris Leininger with noting natural virtual versions of preceding questions, including classification of normal subgroups of finite-index subgroups of \(\operatorname{Mod}_g\), and classification of virtually normal \(H\leq\operatorname{Mod}_g\), defined by
\[
 [\operatorname{Mod}_g:N_{\operatorname{Mod}_g}(H)]<\infty.
\]
These are contextual variations, not a restriction of the main representation problem to finite-index subgroups or to large genus.
\subsection{Short English statement}
For every genus \(g\geq2\) and every graded degree, determine the irreducible symplectic constituents of the Torelli group's Malcev Lie algebra. Hain's reported presentation for \(g>6\) does not itself supply those decompositions; low-genus and representation-category conventions need checking.
\subsection{Sources}
\begin{itemize}
\item Local UnsolvedMath registry, permanent ID 11000015, source code catalogue entry, Problem 2.12, PDF page 28. No individual proposer is supplied.
\item Farb, \emph{Problems on Mapping Class Groups and Related Topics}. Supplied URL: \url{https://www.math.uchicago.edu/~farb/papers/mcgbook.pdf}; recorded access 2026-07-29; extraction pdf-font-structure. The edited volume has 329 explicitly typeset Problem, Conjecture, and Question headings across contributed chapters, with repeated numbering disambiguated by page.
\item Hain [Ha3] is cited for the presentation when \(g>6\); no full bibliographic entry is provided. Chris Leininger is credited for the virtual-version observation in the adjacent discussion.
\item Supplied status open and OPEN-TRIAGE. The summary describes partial results by Hain and others but no complete decomposition of all pieces or complete published answer. These are supplied assessments. Current primary-literature status and redistribution terms for the publicly accessible author/publisher copy remain NEEDS\_REVIEW.
\end{itemize}
\end{document}
